\section{Discussion}
\label{sec:discussion}

\subsection{Geographic Access to Intergovernmental Investment}

After tornado and severe-storm disasters, small counties obtain HMGP projects much less often than larger counties designated in the same declaration, and the shortfall arises because they apply less often, not because their applications fare worse. A grant program whose funding is triggered by hazards can therefore still distribute investment unevenly across places that face similar hazards, because the step that turns an allocation into a project depends on local action. Unequal exposure to hazards is not the source of this disparity. It arises in how the grant is delivered. Post-disaster assistance has similarly been found to widen rather than narrow differences in wealth \parencite{howell2019damages}.

Several explanations can be set aside or bounded. Lower need is unlikely to explain the full gap: the comparisons hold the declaration fixed and control for tornado exposure and damage, and the gap is larger, not smaller, among county--declarations in which a tornado crossed a populated place. Federal or state screening of applications does not explain it, since approval and withdrawal rates are similar once a request is filed. Income, poverty, race, and nonmetropolitan status account for part of the raw difference but leave most of it unexplained, and per-resident fiscal measures add nothing.

A substantial part of the gap is associated with the institutional structure of local government. Small counties have fewer cities, school districts, and special districts that could apply, so they make fewer claims on a declaration's allocation even if each entity behaves like its counterparts elsewhere. This component is a feature of how local government is organized in small counties rather than a failure of any single government. The remaining gap, visible for county governments themselves and in the weaker response to tornado hits, is consistent with the capacity account developed in public administration and planning research \parencite{collins2006redistributive,hall2008grant,lyles2014planners,smith2025rural}: preparing an HMGP application requires expertise and staff time whose cost does not fall in proportion to population. The data do not measure these costs, however, and the residual could also reflect other attributes that vary with county size.

The null result for management-cost reimbursement is informative about timing. If the binding constraint were the ongoing cost of administering a subaward, reimbursing up to 5 percent of it should have helped small counties disproportionately. We find no detectable change, although the estimate is imprecise. This is consistent with a constraint that operates before any award exists, in scoping a project, assembling cost estimates, and completing a benefit--cost analysis, for which reimbursement after an award offers little help. The association with state delivery points in the same direction: where the state prepared more of the applications, the gap is smaller. Neither piece of evidence is causal, but together they suggest that up-front technical assistance is a more promising policy margin than reimbursement after an award.

\subsection{Relation to Previous Findings}

\textcite{seong2022naturally} find that rural HMGP recipients receive less per recipient. Our results suggest that an important margin may be whether there is a recipient at all, and that this margin varies continuously with population rather than with a rural classification. \textcite{miao2025equity} report that aid programs administered through state and local governments favor better-off counties. Our evidence locates one channel for that pattern in the decision to apply, and it shows that, for HMGP after tornado and storm declarations, participation is more closely tied to county size than to income: conditional on size and the other controls, median income is only weakly associated with participation and the poverty rate not at all (Table~\ref{tab:baseline}, column 4; coefficients in Appendix Table~\ref{tab:controls}). \textcite{frimpong2023federal} find that mitigation grants do not raise participation in the Community Rating System. Our findings concern the opposite direction, from local conditions to grant participation, and together the two studies describe a system in which local administrative effort and federal mitigation funding are only loosely connected.

\subsection{Implications for Grant Design}

The findings have three implications for how place-based mitigation grants are designed and evaluated. First, evaluations that compare outcomes across counties with and without mitigation projects should recognize that receiving a project depends on factors associated with county size, which are likely correlated with other determinants of resilience. Second, programs that distribute funds through local applications will tend to direct investment toward larger jurisdictions even when hazard exposure is similar, so equity goals stated in terms of need or vulnerability may not be met without attention to the application stage. Third, the gap is proportionally larger for tornado-specific projects, and small counties in the central and southeastern United States include many places where tornado risk and mobile-home housing are common \parencite{ashley2007fatalities,sutter2010mobile}. Options consistent with the evidence include applications prepared by the state or by regional bodies on behalf of small jurisdictions, pre-application technical assistance, and simplified cost-effectiveness requirements for standard project types. The present analysis does not evaluate the effects of these options.

\subsection{Limitations}
\label{sec:limitations}

County size is not randomly assigned, and the residual gap could reflect unmeasured attributes correlated with size, such as the presence of an emergency manager, local political preferences regarding federal programs, or the stock of buildings suitable for retrofit. The capacity proxies available in national data are coarse; per-resident spending measures resources, not expertise. The decomposition by number of local governments is observational and does not identify the effect of changing local government structure.

Several measurement choices also limit the analysis. Tornado exposure is assigned by the calendar year in which the declared incident began, not by the incident period itself, so a hit may reflect a tornado from a different event in the same year. The subapplication records linked to counties contain few older closed or denied requests, so the comparison of decision rates relies mainly on recent records and on the approved, withdrawn, and not-approved categories. Requests for declarations late in the sample may still be pending.

The within-county and reform tests have limited power, and the reform coincided with other changes in federal mitigation policy. The analysis covers HMGP after tornado and severe-storm declarations and may not extend to flood programs, competitive pre-disaster grants, or other countries. Finally, the outcomes are applications and awards, not completed projects or avoided losses, so the results concern access to funding rather than differences in realized risk reduction.
